\chapter*{Preface}
\addcontentsline{toc}{chapter}{Preface}

These notes accompany \emph{Models of Neural Systems}. The course moves from linear classifiers and learning rules to membrane biophysics, excitability, cable theory, network dynamics, and stochastic spiking. The unifying question is always the same: which state variables are sufficient for the phenomenon we want to explain?

The emphasis is mathematical, but not formal for its own sake. Derivations are included when they expose mechanism: why diffusion and drift set a reversal potential, why voltage-dependent conductances create spikes, why bifurcations change coding, and why morphology reshapes propagation. At each stage, the model is only as detailed as the question demands.

Throughout, we return to three habits of thought:
\begin{itemize}
  \item \textbf{Choose the right variables.} Channel states, membrane voltage, spike times, firing rates, and latent population coordinates are different coarse-grainings of the same system.
  \item \textbf{Track the parameters that control behavior.} Time constants, gain terms, eigenvalues, and length constants are more informative than isolated formulas.
  \item \textbf{Read equations qualitatively.} Before solving, ask what each term pushes, damps, couples, or conserves.
\end{itemize}

The notes draw on classic references in theoretical and computational neuroscience \parencite{dayan2001theoretical,gerstner2014neuronal,koch1999biophysics,izhikevich2007dynamical,hodgkin1952quantitative}. Read them actively: identify the state variables, limiting cases, and qualitative predictions before memorizing the algebra.

\bigskip
\noindent\textit{Berlin, February 2026}
